\scanpage{60}
% archive sheet number: 59
\uline{Schéma d'Albanese}.

Soit $X$ \deleted{projectif et} \inserted{fid} plat \inserted{\unsure{q. cpt.|qu.-proj.}} sur $S$, \deleted{$\underline O_S\simeq f_*(\underline O_X)$ univt}

\deleted{On dit} Données équivalentes

\deleted{Pour un espace}

(i) $P$ fibré principal homogène [pour descente fid plate et \unsure{q. cpte|qu.-proj.}]
sur $S$, $S$-groupe $G$, et un $S$-morphisme $\unsure{\psi|u|\varphi}:X\longrightarrow P$

(\deleted{\unsure{ii}}i bis) Un $S$-groupe $G$, et une donnée de descente effective sur
$G_X=G\times_S X$ pour $X\longrightarrow S$, et la structure de \deleted{groupe} fibré
pr. hom. \inserted{à droite} sous $G_X$ sur $G_X$.

% left margin, next to (i bis)
\diagram{\begin{tikzcd}[ampersand replacement=\&,row sep=small,column sep=small]
P\arrow[d]\&P_X\arrow[l]\arrow[d]\\
S\&X\arrow[l]\arrow[ul,dotted,"\psi"']
\end{tikzcd}}

(ii) Un \deleted{\unsure{hom}} $S$-groupe $G$, et un $S$-hom $\varphi:X\times_S X$ satisfaisant
la condition \inserted{symbolique} $\varphi(x,y)\varphi(y,z)=\varphi(x,z)$ [impliquant $\varphi(x,x)=e$],
\deleted{\unsure{et}} et une condition d'effectivité, [autom. vérifiée
si $G$ est affine sur $S$, ou si $X$ est \unsure{qu. pr.} sur $S$
et $X\longrightarrow S$ admet loc. des quasi-sections plates finies,
p. ex. $X\longrightarrow S$ projectif].

% left margin, beside the cocycle condition
\diagram{\begin{tikzcd}[ampersand replacement=\&,column sep=small]
S \& X \& X\times_S X\arrow[l,shift left=1.3]\arrow[l,shift right=1.3] \& X\times_S X\times_S X\arrow[l,shift left=2]\arrow[l]\arrow[l,shift right=2]
\end{tikzcd}}

Supposons enfin que \deleted{\unsure{$G$}} $X$ soit \deleted{\unsure{projectif}} propre sur $S$
et tel que $\underline O_S\simeq f_*(\underline O_X)$ univt, alors on a équivalence
avec

(iii) Un $S$-groupe $G$, et un $S$-hom $\varphi:X\times_S X$ satisfaisant
seulement $\varphi(x,x)=e$, et une condition d'effectivité
[qui est autom. vérifiée si $G$ est affine/$S$, ou $G$ :
\unsure{qu. proj.}/$S$ et $X/\unsure{s|S|e}$ admet loc. des quasi-sections finies
plates, p. ex. $X$ projectif sur $S$].

% slanted
\margin{définir aussi en termes de $G$, $\unsure{\mathrm{Hom}|\mathrm{hom}}_{X/S}(G)={}$ ens des $(P,u)$ avec $G$ donné, d'où foncteur \deleted{\unsure{et}} et \unsure{faisceau|préfaisceau} de Weil, on peut dire [\unsure{par} ``\unsure{préfaisceau} de Weil''] $S'\longmapsto\underbrace{\mathrm{Hom}_{S'}(X_{S'},G_{S'})}_{\unsure{h^0|h}(X_{S'}/S',G_{S'})}/\unsure{\mathrm{Hom}^{\mathrm{cts}}|\mathrm{Hom}^{\mathrm{ct}}}$ et on passe au \uline{faisceau} associé \unsure{---|\ldots}}

Pour simplifier, nous supposons par la suite
$f:X\longrightarrow S$ \uline{projectif}, et $\underline O_S\simeq f_*(\underline O_X)$ univt. Pour
fibré principal sous variable $P$, \inserted{\unsure{ou|au} groupe variable $G$ \deleted{\unsure{propre}} \inserted{\unsure{comm.|proj.}} \unsure{proj.} sur $S$}, considérons
$\mathrm{Hom}_S(X,P)$ comme foncteur covariant en $P$, \inserted{soit $T_X(P)$.} On veut
le représenter. Tout d'abord, ce foncteur est \uline{pro-représentable}. Comme dans la catégorie des $P$, les \uline{lim}
